% F05：本书原创矢量示意；依据所在节的定义、证明或明确模型。
\begin{figure}[htbp]
\centering
\begin{tikzpicture}[x=1cm,y=1cm,>=stealth,every node/.style={font=\small},line width=.65pt]

\fill[AnalysisBlue!12] (2,1.5) ellipse (1.6 and 1.1);
\draw[AnalysisBlue,thick] (2,1.5) ellipse (1.6 and 1.1);
\node at (2,1.5){\(C\)};
\draw[orange!80!black,thick] (4.3,-.2)--(4.3,3.2);
\node[above] at (4.3,3.2){\(f=\alpha\)};
\fill (6,1.6) circle (2pt) node[right]{\(x_0\)};
\draw[->] (4.6,.4)--(6.4,.4) node[midway,below]{\(f\) 增大};
\fill[AnalysisBlue!12] (9.5,1.5) circle (1.1);
\draw[AnalysisBlue,thick] (9.5,1.5) circle (1.1);
\draw[orange!80!black,thick] (10.6,-.2)--(10.6,3.2);
\fill (10.6,1.5) circle (2pt);
\node at (9.5,-.8){支撑：边界允许相交};
\node at (2,-.8){分离：\(\sup_C f<f(x_0)\)};

\end{tikzpicture}
\caption[凸分离与支撑泛函]{凸分离与支撑泛函。左图示意实泛函的严格分离；右图示意支撑超平面。二维欧氏图只解释几何位置，不要求一般分离泛函来自正交投影。}
\label{b1:fig:visual-f05}
\end{figure}
